\section{Methods}
\label{sec:methods}

Every parameter, with its source, range and whether it is measured, derived or assumed, is listed in
Table~\ref{SI-tab:params}. Code and data are archived (Section~\ref{SI-sec:si-data}).

\subsection{Source}
\label{sec:source}
Passage volumes span 17--375\,mL, with a median of 92\,mL \cite{suarez1997}. In the Monte Carlo analysis, volume $V$,
duration $T$ (0.3--20\,s) and effective orifice diameter $d$ (1--8\,mm) are drawn log-uniformly. Neither has been
measured for a passage. The upper bounds come from anorectal physiology: flatus passes during an anal relaxation
lasting $19 \pm 6$\,s \cite{ronholt1999,corsetti2017}, and the mid anal canal widens to about 8\,mm only at balloon
pressures above those that drive flatus \cite{luft2012,sorensen2014,miller1988b}. The lower bounds are assumed.

Sulfur gas concentrations are taken from 87 passages of 16 volunteers \cite{suarez1998}: \ce{H2S}
1.06\,\textmu mol\,L$^{-1}$, methanethiol 0.21\,\textmu mol\,L$^{-1}$ and dimethyl sulfide 0.08\,\textmu mol\,L$^{-1}$.
They vary between passages by a log-normal factor with $\sigma_{\ln} = 1.1$, derived from the reported standard
errors, drawn independently of volume because sulfur concentration does not correlate with passage volume
\cite{suarez1997}. Buoyancy is computed from measured compositions and is signed (Section~\ref{sec:methods-chem}); it
is drawn independently of volume because, across the 16 subjects of ref.~\citenum{suarez1997}, it does not correlate
with mean passage volume (Spearman $-0.15$).

\subsection{Odorant chemistry}
\label{sec:methods-chem}
\emph{Odour activity.} Each sulfur gas contributes an odour activity value (OAV), its mole fraction over its detection
threshold, $\mathrm{OAV}_i = x_i/\theta_i$. The passage's OAV is their sum; we assume additivity and do not model
mixture suppression. We compare two measured compositions:
\begin{itemize}
  \item 87 passages of 16 volunteers \cite{suarez1998}, used throughout the dispersion model;
  \item 28 passages measured by gas chromatography: \ce{H2S} 149, methanethiol 216 and dimethyl sulfide
  144\,nmol\,L$^{-1}$ \cite{tangerman2009}.
\end{itemize}
Three threshold sets bracket the observer: a trained panel \cite{nagata2003}, population geometric means
\cite{amoore1983,wilby1969} and the lowest values in the literature, as compiled in ref.~\citenum{tangerman2009}.

\emph{Removal check.} Zinc acetate precipitates \ce{H2S} and thiols; activated charcoal adsorbs all three gases. For both
treatments, ref.~\citenum{suarez1998} reports the gas concentrations and two judges' ratings before and after. We
predict the rating ratio as $(\mathrm{OAV}_\mathrm{after}/\mathrm{OAV}_\mathrm{before})^n$ with $n = 0.20$--$0.45$
\cite{collins2001}, without fitting.

\emph{Buoyancy.} The reduced gravity $g' = g(\rho_a-\rho_g)/\rho_a$ of each of 16 measured compositions
\cite{suarez1997} is split into a thermal part, $g(1-T_a/T_g)$, for gas at 37\,\textdegree C in air at
22\,\textdegree C, and a compositional part, $g(M_a-M_g)/M_a$, at equal temperature with air at 50\% relative
humidity. The water-vapour content of flatus has not been measured, so we bracket dry and saturated gas.

\emph{Gas-phase and surface losses.} Gas-phase losses use IUPAC rate constants at 298\,K for reaction with OH,
\ce{O3}, \ce{NO3} and Cl \cite{atkinson2004}, with oxidant levels measured in an office: OH
$0.6$--$9.8\times10^{5}$\,cm$^{-3}$ and \ce{O3} 0.7--2.6\,ppb \cite{fiorentino2023}. Two assumed upper cases add OH
$10^6$\,cm$^{-3}$ with \ce{O3} 40\,ppb, and Cl $10^6$\,cm$^{-3}$ during hypochlorite cleaning. No deposition velocity
$v_d$ has been published for these gases on indoor surfaces, so surface loss is bracketed as $k = v_d\,S/V$ with
$S/V = 3$\,m$^{-1}$.

\emph{Dilution in a panel test.} Ref.~\citenum{tangerman2009} reports the lowest source concentration that every panel
member detected when 100\,mL was released below the nose. Divided by a detection threshold, it gives the dilution
between release and nose in that test, which we compare with the breath-averaged dilution of the analytic model.

\subsection{Acoustics}
\label{sec:acoustics}
The anal canal is modelled as a lip reed, a soft-tissue valve that self-oscillates in the flow through it, as the
lips do on a brass instrument. It oscillates when the driving pressure $\tfrac12\rho(u_0/C_\mathrm{d})^2$ exceeds a
threshold $P_\mathrm{th}$ = 0.6--2.5\,kPa (nominal 1.2\,kPa), with discharge coefficient $C_\mathrm{d} = 0.85$. The
closest measured soft-tissue sphincter source, the pharyngo-oesophageal segment of oesophageal speech, sustains
phonation down to 1.1--2.5\,kPa \cite{nagai1993,grolman2007,ng2011}, which bounds the onset from above; lips on a brass
instrument start at about 0.6\,kPa \cite{matteoli2021}. This gives an onset velocity
$u_\mathrm{onset} \approx 42$\,m\,s$^{-1}$ (21--71).

An oscillating valve radiates as a harmonic monopole with harmonics falling by 12\,dB per octave. Its fundamental is
taken as 160\,Hz (60--300\,Hz), between oesophageal speech at 60--150\,Hz \cite{nagai1993,robbins1984,pindzola1989,deschler1999}
and recordings of flatus near 270\,Hz, which are not peer-reviewed. A release counts as \emph{audible} when at least one harmonic exceeds both the ISO 226 hearing threshold
\cite{iso226_2023} and the masked threshold of the background, computed from the critical ratio \cite{hawkins1950}.
Backgrounds range from a 35\,dB(A) living room to a 77\,dB(A) subway car.

\subsection{Analytic dispersion model}
\label{sec:gate0}
Dispersion is modelled in two phases.
\begin{enumerate}
  \item \emph{Near field.} The release behaves as a starting jet. Its tip follows
  $x = \Gamma (\dot Q_0 u_0)^{1/4} t^{1/2}$, and once the source stops it follows the puff law
  $x^4 = C_p^4 Q_0 u_0 t$, with $\Gamma = 2.8$ and $C_p = 2.6$ \cite{sangras2002,diez2003}. The centreline dilution
  is $K_c d^*/x$ with $K_c = 5.11$. In parallel, a buoyant thermal grows with entrainment coefficient
  $\alpha = 0.25$ \cite{landeau2014}, and the larger of the two dilutions is kept. This phase ends when the
  self-induced velocity has fallen to the ambient velocity $u_h$ (0.05--0.25\,m\,s$^{-1}$), which gives
  Eq.~\eqref{eq:silencer}.
  \item \emph{Transport.} The cloud is then advected to the nose and spreads as a Gaussian puff,
  $\sigma^2 = \sigma_0^2 + 2Kt$, with indoor eddy diffusivity $K = 10^{-4}$--$1.3\times10^{-2}$\,m$^2$\,s$^{-1}$
  \cite{cheng2011}.
\end{enumerate}
Two receptors are modelled: the releaser's own nose, reached by the convective boundary layer, the layer of warm air
rising along the body \cite{licina2014}, at a speed taken from the CFD trajectories below; and a neighbour
0.8--1.4\,m away, reached by room air. Concentrations are averaged over one inspiration.

We draw 20\,000 pairs of silent and audible releases, each classified by the acoustic onset and evaluated with
environmental parameters drawn from their ranges. Three simplifications favour the silent release: turbulent-jet
constants are used even for laminar silent releases, the nose sits on the cloud centreline, and the centreline
concentration is used.

\subsection{Computational fluid dynamics}
\label{sec:cfd}
We use OpenFOAM 13 with transient Reynolds-averaged Navier--Stokes equations (RANS), the $k$--$\omega$ SST turbulence
model, the Boussinesq approximation for buoyancy and the energy equation in internal-energy form.

\emph{Validation} (Section~\ref{SI-sec:si-validation}). A free round jet at $\mathrm{Re} = 6700$ reproduces the
canonical decay constants (velocity 5.0, scalar $K_c = 4.3$--4.6) with a grid uncertainty of at most 3.7\%. The
convective boundary layer of a seated manikin is 40--56\% faster than thermal-manikin data \cite{licina2014} and
follows their trend with room temperature. Both biases favour the silent release: excess jet spreading over-dilutes
audible releases by about 15\%, and the fast boundary layer shortens travel times to the releaser's own nose by about
30\%, which the analytic model corrects.

\emph{Scenes} (Section~\ref{SI-sec:si-scenes}). Manikins built from geometric primitives, with skin at
33\,\textdegree C, stand or sit in rooms at 23\,\textdegree C:
\begin{itemize}
  \item a sofa: two seated people 0.6\,m apart in a 3.2 $\times$ 3.2 $\times$ 2.6\,m living room, with the releaser
  upright, reclined by 4.5\,\textdegree\ or leaning forward by 10\,\textdegree, over 30\,s (upright also 120\,s);
  \item an elevator: four standing people in a closed 630\,kg car, fan off, over 60\,s, with the releaser in the front
  or the back row;
  \item a standing pair: two people face to face, noses 1.0\,m apart, in the same living room, in still air and with
  office ceiling ventilation at 2.9 air changes per hour (EN 16798-1 Category II \cite{olesen2024}), over 60\,s.
\end{itemize}
Each scene is first run without a release until the body plumes have developed. From that flow field, the same
92\,mL is released behind the releaser's buttocks, as an inlet that matches the release's volume flow and momentum:
silent, 2\,s through 8\,mm at $u_0 = 0.9$\,m\,s$^{-1}$; audible, 0.5\,s through 2.4\,mm at 41\,m\,s$^{-1}$. The
odorant is a passive scalar, sampled 2.5\,cm in front of each nose, over 5\,cm discs at nose height and over
sinusoidal breathing cycles.

\subsection{Perception and observers}
\label{sec:perception}
The nose samples only during inhalation: sinusoidal flow at 12 breaths per minute, inspiratory fraction 0.42
\cite{wysocki2006}. Perceived intensity follows Stevens' power law, $I \propto \mathrm{OAV}^n$, with $n = 0.28$
measured for flatus \ce{H2S} \cite{collins2001}. Two passages are distinguishable when their concentrations differ by
more than the olfactory Weber fraction of 0.28 \cite{stone1965}, i.e.\ a just-noticeable ratio of 1.28.

The OAVs of the dispersion model use trained-panel thresholds \cite{nagata2003}. For untrained observers:
\begin{itemize}
  \item population thresholds are 10--23 times higher \cite{amoore1983,wilby1969};
  \item detection is a per-passage probability $\Phi(\log_2\mathrm{OAV}/\sigma)$ with $\sigma = 2$ binary steps
  \cite{amoore1983}, since one person's threshold varies over time about as much as thresholds vary across people
  \cite{stevens1988};
  \item a bystander who is not expecting an odour needs 26 times the threshold for a 50\% response, with
  $\sigma = 3.5$ binary steps \cite{amoore1983,whisman1978}.
\end{itemize}
Distraction reduces noticing \cite{forster2018} but not the peak rated intensity \cite{fallon2018}, so intensity is
computed from the attentive OAV. Two variants follow: clothing removes 22\% of the sulfide gas \cite{ohge2005}, and a
sound can make a distracted bystander attentive for that passage.

\subsection{Uncertainty analysis}
\label{sec:uq}
Uncertainty is propagated through the analytic model; one CFD run of 30--60\,s costs 1--3\,h on four cores, too many
for a sampling design. First-order Sobol indices (the share of the variance of $\log_{10} R$ explained by one input
alone) are estimated by binning the Monte Carlo sample, with bootstrap 90\% intervals. Because the silent and audible
releases of a pair are drawn by rejection, the inputs are dependent and the indices are descriptive. The regime map
shows $\log_{10} R$ at equal mass against the audible release's Silencer number and the travel time.

\emph{Passages as emitted in the simulated scenes.} Each CFD scene gives the doses of one release volume at equal mass.
The OAV is linear in concentration, so a passage of volume $V$ following the same flow path gives
$\mathrm{OAV}\times V/92$\,mL. For each of the 20\,000 Monte Carlo pairs we scale the silent and the audible CFD doses
by the pair's two volumes, and report $P(R>1)$ and the shares above and below the just-noticeable ratio per scene and
nose. In the analytic model this scaling is exact ($R_\text{as emitted} = R_\text{equal mass}\,V_s/V_l$). Applied to
the CFD it assumes that the flow path does not depend on the volume: close for silent releases, whose momentum is
negligible against the body plume, and an approximation for audible ones.
